\scanpage{60}
% archive foliation: 59

Alors pour tout objet $Y$ de $\unsure{X|C}$, \deleted{$i(X)(Y)$} posant $F=i(X)$,
$F(Y)$ est un module \inserted{à droite} projectif de type fini sur $O'{}^o$, où $O'=\operatorname{End}(Y)$.
Soit \deleted{\unsure{$C_f$}}, dans le cas général,
\deleted{\unsure{$\mathfrak C$}}
$C_{\mathrm{sp}}^{\unsure{i}}(A)$ la catégorie de ces foncteurs \inserted{(dits ``spéciaux'')}.
On trouve donc ici une immersion (sous réserve d'existence des gén. stricts)
\margin{\deleted{\unsure{en toute} \unsure{espèce de type} ``\unsure{fini}'' \unsure{en sens général}}}
\diagram{\begin{tikzcd}
C \arrow[r,"i"] & C_{\mathrm{sp}}(A)
\end{tikzcd}},
et le choix d'un générateur \inserted{strict} $X$ de $C$ permet d'identifier
$C_{\mathrm{sp}}(A)$ à la catégorie des $O$-modules à droite projectifs de type fini.
\margin{\unsure{il peut être préférable de se placer}: \unsure{loc.} libre de rg fini sur $\operatorname{Spec}A=S$.}
\deleted{\unsure{La}}. On l'appelle la
\deleted{\unsure{enveloppe projective}}
\deleted{\inserted{\unsure{saturé projectif}}}
(ou le \inserted{\unsure{simple}} \unsure{saturé}) de $C$.
